\documentclass[11pt,a4paper]{article}
\usepackage[margin=1.8cm]{geometry}
\usepackage{graphicx}\usepackage{amsmath,amssymb,tabularx,array,needspace,fontspec,enumitem,cancel}
\newcommand{\tb}{\cancel{||||}\,}
\setmainfont{FreeSerif}
\newfontfamily\urdufont[Script=Arabic]{FreeSerif}
\TeXXeTstate=1
\newcommand{\urdu}[1]{{\urdufont\beginR #1\endR}}
\hyphenpenalty=10000\exhyphenpenalty=10000
\setlength{\parindent}{0pt}
\setlength{\parskip}{0pt}
\renewcommand{\arraystretch}{1.15}
\pagestyle{plain}

% \q{roman}{text}
\newlength{\qw}\settowidth{\qw}{\textbf{CCCLXXXVIII.}\ }
\newcommand{\q}[2]{%
  \Needspace{3\baselineskip}%
  \par\medskip\noindent\hangindent=\qw\hangafter=1\makebox[\qw][l]{\textbf{#1.}}#2\par}

\begin{document}

\begin{center}
  {\Huge\bfseries The Hope Science Academy}\\[3pt]
  {\small Excellence is our standard not goal}\\[2pt]
  {\small 247-E-1, (Opposite LDA School) Johar Town, Lahore. PH: 0300-8194789, 0322-4157001}
\end{center}

\medskip
\noindent
\begin{tabularx}{\linewidth}{|>{\raggedright\arraybackslash}X|>{\raggedright\arraybackslash}X|>{\raggedright\arraybackslash}X|>{\raggedright\arraybackslash}X|}
\hline
\textbf{Student Name:} & \textbf{Roll Num:} & \textbf{Class:} 12th & \textbf{Paper Code:} \\ \hline
\textbf{Subject Name:} Mathematics. & \textbf{Time Allowed:} & \textbf{Total Marks:} 374 & \textbf{Exam Date:} 03-Nov-2026 \\ \hline
\multicolumn{4}{|l|}{\textbf{Exam Syllabus:} FULL BOOK} \\ \hline
\end{tabularx}

\bigskip
\noindent\textbf{Q1. Write short answers of the following questions.}\hfill\textbf{1$\times$374=374}

\q{I}{Sketch the graph of the following function using the method of factorization: y = 3(x+1) (x-1)}
\q{II}{Sketch the graph of the following function using the method of factorization: f(x) = -2(x-1)(x-2)}
\q{III}{Sketch the graph of the following function using the method of factorization: y =(2x - 1)(2x + 1)}
\q{IV}{Sketch the graph of the following function using the method of factorization: y = 2x$^{2}$ + x-3}
\q{V}{Sketch the graph of the following function using the method of factorization: f(x) = x$^{2}$ + 2x + 1}
\q{VI}{Sketch the graph of the following function using the method of factorization: y = 3(x - 1)$^{2}$}
\q{VII}{Find the function whose graph is shown in the following figure:\par\smallskip\includegraphics[height=3.6cm]{figA.png}}
\q{VIII}{Find the function whose graph is shown in the following figure:\par\smallskip\includegraphics[height=3.6cm]{figB.png}}
\q{IX}{Find an equation of the graph of a parabola whose x-intercepts are x = -1 and x = -2 and which passes through the point (-3, 2).}
\q{X}{Find an equation of the graph of a parabola which touches x-axis at x=-1 and cuts y-axis at y=-2.}
\q{XI}{Find an equation of a quadratic function whose graph cuts x-axis at x = -1 and x = 1, and whose vertex is at (0, -2).}
\q{XII}{Sketch the graph of y=2(x-1)(x-3).}
\q{XIII}{Sketch the graph of y = -4x$^{2}$ - 8x - 4}
\q{XIV}{Find the function whose graph is shown in Figure.\par\smallskip\includegraphics[height=3.6cm]{figC.png}}
\q{XV}{Find an equation of the graph of a parabola whose x-intercepts are x = 1 and x = 2 and passes through the point (3, 4).}
\q{XVI}{Is the following function is algebraic? $f(x)=\sec x$}
\q{XVII}{Is the following function is algebraic? $g(x)=\dfrac{2x-3}{x+1}$}
\q{XVIII}{Is the following function is algebraic? $h(x)=\sqrt{x^{2}+2x}$}
\q{XIX}{Is the following function is algebraic? $k(x)=\dfrac{e^{x}-e^{-x}}{2}$}
\q{XX}{Is the following function is algebraic? $r(x)=(x-1)^{3}$}
\q{XXI}{Is the following function is transcendental? $f(x)=\dfrac{1}{x}$}
\q{XXII}{Is the following function is transcendental? $g(x)=\cot x$}
\q{XXIII}{Is the following function is transcendental? $h(x)=e^{2x}+\sqrt{x^{2}+2x}$}
\q{XXIV}{Is the following function is transcendental? $k(x)=\dfrac{e^{x}+e^{-x}}{2}$}
\q{XXV}{Is the following function is transcendental? $r(x)=(x-1)^{-\frac{3}{2}}$}
\q{XXVI}{Prove the following hyperbolic identities: $1-\tanh^{2}x=\mathrm{sech}^{2}x$}
\q{XXVII}{Prove the following hyperbolic identities: $\mathrm{csch}^{2}x=\coth^{2}x-1$}
\q{XXVIII}{Prove the following hyperbolic identities: $\sinh 2x=2\sinh x\cosh x$}
\q{XXIX}{Prove the following hyperbolic identities: $\cosh 2x=\cosh^{2}x+\sinh^{2}x$}
\q{XXX}{Prove that $\cosh^{2}x-\sinh^{2}x=1$}
\q{XXXI}{Given the function $f(x)=\sqrt{x}$ find the equation of the function after: vertical stretch by factor 2.}
\q{XXXII}{Given the function $f(x)=\sqrt{x}$ find the equation of the function after: horizontal shift 3 units to the right.}
\q{XXXIII}{Let $f(x)=\ln(x^{2}+1)$. Find the function obtained after: a vertical stretch by factor 2.}
\q{XXXIV}{Let $f(x)=\ln(x^{2}+1)$. Find the function obtained after: then a shift 1 unit to the left.}
\q{XXXV}{Solve the following equation: $\log(3x-2)=\log(x^{2}-5x+6)$}
\q{XXXVI}{Solve the following equation: $\ln(x^{2}-1)=\ln(3x-2)$}
\q{XXXVII}{Solve the following equation: $x\,\log_{\frac{1}{e}}4=\log_{\frac{1}{e}}(16^{x}-3)$}
\q{XXXVIII}{Solve the following equation: $2^{x}+2^{1-x}\leq 3$}
\q{XXXIX}{Solve the following inequalitie (if its solution exists): $2^{x}\geq\log_{2}(16)$}
\q{XL}{Solve the following inequalitie (if its solution exists): $\log_{2}(x-1)\geq\log_{2}(5-x)$}
\q{XLI}{Solve the equation: $4^{x-1}=8$}
\q{XLII}{Solve the inequality: $\log\,(2x-1)\geq\log(5-3x)$}
\q{XLIII}{A population is given by N = 300e$^{0.2t}$ , where t is measured in hours. Find the time when the population becomes 600.}
\q{XLIV}{A culture grows according to $N=150e^{kt}$ , where time t is measured in hours. After 10 hours, the population becomes 250. Find the growth constant k.}
\q{XLV}{A radioactive substance decays as $N=800e^{-0.05t}$ , where t is in hours. Find the time when the amount redues to 200.}
\q{XLVI}{A substance decays according to $N=600e^{kt}$ , where time t is measured in hours, After 8 hours, the amount becomes 400. Find the decay constant k.}
\q{XLVII}{A sound has intensity $I=10^{-6}W/m^{2}$ and reference intensity $I_{o}=10^{-12}W/m^{2}$ Find the sound level in decibels.}
\q{XLVIII}{A sound has level L=80 dB and $I_{o}=10^{-12}W/m^{2}$ . Find the intensity I.}
\q{XLIX}{Two sounds have intensities $I_{1}$ and $I_{2}$ such that $L_{1}=85$ and $L_{2}65$ dB . Find the ratio $\dfrac{I_{1}}{I_{2}}$ where $I_{o}=10^{-12}W/m^{2}$ .}
\q{L}{A sum of Rs. 4000 is invested at 6\% per annum, compounded annually for 3 years. Find the final amount.}
\q{LI}{A population is given by $N=200e^{0.04t}$: find the time t when the population becomes 400, where t is measured in hours.}
\q{LII}{A radioactive substance decays according to : $N=500e^{kt}$ After 20 hours the amount becomes 300. Find the value of k. where t is in hours.}
\q{LIII}{A sound has intensity $I=10^{-4}$ and reference intensity $I_{o}=10^{-12}$ . Find the sound level.}
\q{LIV}{A sound has level L=60dB and $I_{o}=10^{-12}$ . Find the intensity I.}
\q{LV}{Differentiate with respect to the variable involved: $y=x^{2}\cos$}
\q{LVI}{Differentiate with respect to the variable involved: $y=x\sec x+\dfrac{1}{x}$}
\q{LVII}{Differentiate with respect to the variable involved: $y=\dfrac{\sin t}{1-\cos t}$}
\q{LVIII}{Differentiate with respect to the variable involved: $y=(1+\csc t)\sin t$}
\q{LIX}{Differentiate with respect to the variable involved: $y=\sec x+4\sqrt{x}-10$}
\q{LX}{Differentiate w.r.t the variable involved: $y=\sec(\tan x)$}
\q{LXI}{Differentiate w.r.t the variable involved: $y=\sqrt{\sin x}+\cos\sqrt{x}$}
\q{LXII}{Differentiate w.r.t the variable involved: $y=\tan^{2}x+\cot(x^{2})$}
\q{LXIII}{Differentiate w.r.t the variable involved: $y=\sin^{3}(\cos 2v)$}
\q{LXIV}{Find $\dfrac{dy}{dx}$ for the following parametric equation: $x=\sin^{2}t$ and $y=\cos 2t$}
\q{LXV}{Use implicit differentiation to find $\dfrac{dy}{dx}$ if: $x+\tan(xy)=0$}
\q{LXVI}{Differentiate w.r.t the variables involved: $\sin^{-1}\!\left(\dfrac{x}{a}\right)$}
\q{LXVII}{Find $f'(x)$ if $f(x)=\tan x$}
\q{LXVIII}{Find $\dfrac{dy}{dx}$ if $y=x^{2}\csc x$}
\q{LXIX}{Differentiate $y=\sin(x^{2}-x)$ w.r.t 'x'.}
\q{LXX}{Find $\dfrac{dy}{dx}$ if $x+y^{2}+\sin(xy^{2})=0$.}
\q{LXXI}{Find $f'(x)$ if: $f(x)=(x^{3}+4x)e^{1/x}$}
\q{LXXII}{Find $f'(x)$ if: $f(x)=\sin(2^{x})$}
\q{LXXIII}{Find $f'(x)$ if: $f(x)=e^{x\cos x}$}
\q{LXXIV}{Find $f'(x)$ if: $f(x)=\dfrac{e^{x}+e^{-x}}{e^{x}-e^{-x}}$}
\q{LXXV}{Find $\dfrac{dy}{dx}$ if: $y=\ln(e^{x}-e^{-x})$}
\q{LXXVI}{Find $\dfrac{dy}{dx}$ if: $y=\ln\!\left(x-\sqrt{x^{2}-1}\right)$}
\q{LXXVII}{Find $\dfrac{dy}{dx}$ if: $y=\ln(1-\cos^{2}x)$}
\q{LXXVIII}{Find $\dfrac{dy}{dx}$ if: $y=\dfrac{\ln x}{x^{2}}$}
\q{LXXIX}{Find $\dfrac{dy}{dx}$ if $y=e^{\sin^{2}x}$.}
\q{LXXX}{Find $\dfrac{dy}{dx}$ if $a^{x^{3}}$.}
\q{LXXXI}{Find $\dfrac{dy}{dx}$ if $y=\dfrac{a^{x}}{\log_{a}x}$ .}
\q{LXXXII}{Find $\dfrac{dy}{dx}$, if $y=\ln(\cos(e^{x}))$.}
\q{LXXXIII}{Find the equation of the tangent line to the curve $y=\sin^{-1}\!\left(\dfrac{x}{2}\right)+\cos x$ at the point where $x=\dfrac{\pi}{2}$}
\q{LXXXIV}{Find $y_{2}$ if: $y=x^{4}-3x^{2}+4x-5$}
\q{LXXXV}{Find $y_{2}$ if: $y=\ln(2x+3)$}
\q{LXXXVI}{Find $y_{2}$ if: $y=x^{3}e^{2x}$}
\q{LXXXVII}{Find $y_{2}$ if: $x^{2}-y^{2}=4$}
\q{LXXXVIII}{Find $y_{2}$ if: $x^{2}-xy+y^{2}=7$}
\q{LXXXIX}{Prove that $f(x)=e^{x}$ defined on $R=(-\infty,\infty)$ is increasing/}
\q{XC}{Show that $f(x)=2^{-x}$ defined on $R=(-\infty,\infty)$ is decreasing.}
\q{XCI}{Find the differential dy in the following case in the following case: $y=x^{3}-3\sqrt{x}$}
\q{XCII}{Find the differential dy in the following case in the following case: $y=\cos x^{2}$}
\q{XCIII}{Use dy to approximate $\triangle y$ for the following function: $y=x^{3}-x$ when x changes from 3 to 3.02.}
\q{XCIV}{Find dy if y=x4 + 7x.}
\q{XCV}{Find the value of dy when x=1 and dx=0.1.}
\q{XCVI}{Find dy and $\triangle y$ for the function $f(x)=x^{3}$ when x=1 and dx=0.1.}
\q{XCVII}{Find all the antiderivatives for the following derivative: $\dfrac{dy}{dx}=100x^{4}$}
\q{XCVIII}{Find all the antiderivatives for the following derivative: $\dfrac{dy}{dt}=2+24t^{3}-10t^{4}$}
\q{XCIX}{Find indefinite integral, (check answer by differentiating): $\displaystyle\int(1+2t)\,dt$}
\q{C}{Find indefinite integral, (check answer by differentiating): $\displaystyle\int(3u^{2}+4u)\,du$}
\q{CI}{Find indefinite integral, (check answer by differentiating): $\displaystyle\int100x^{8}\,dx$}
\q{CII}{Find indefinite integral, (check answer by differentiating): $\displaystyle\int(2x^{-2}+3x^{-3})\,dx$}
\q{CIII}{Evaluate the following: $\displaystyle\int(3x+3^{x}+x^{3})\,dx$}
\q{CIV}{Evaluate the following: $\displaystyle\int x^{2}\!\left(x+\dfrac{1}{x}\right)dx$}
\q{CV}{Evaluate the following: $\displaystyle\int\left(\sqrt{y}+\dfrac{1}{y}\right)^{2}dy$}
\q{CVI}{Evaluate the following: $\displaystyle\int\dfrac{\sqrt{x}\,(x+2)}{x}\,dx$}
\q{CVII}{Evaluate the following: $\displaystyle\int\dfrac{(e^{x}+e^{-x})}{e^{3x}}\,dx$}
\q{CVIII}{Evaluate the following: $\displaystyle\int7^{\ln}\,dx$}
\q{CIX}{Evaluate the following: $\displaystyle\int2^{\log_{4}x}\,dx$}
\q{CX}{Evaluate the following: $\displaystyle\int\csc x\,dx$}
\q{CXI}{Evaluate the following: $\displaystyle\int\cot^{2}x\,dx$}
\q{CXII}{Evaluate the following: $\displaystyle\int\sin x\,.\,\cos 7x\,dx$}
\q{CXIII}{Evaluate the following integrals: $\displaystyle\int x^{99}\,dx$}
\q{CXIV}{Evaluate the following integrals: $\displaystyle\int-98x^{-99}\,dx$}
\q{CXV}{Evaluate the integral $\displaystyle\int(8x^{7}-6x^{5}+4x^{3}-2x-3)\,dx$}
\q{CXVI}{Evaluate the integrals: $\displaystyle\int\dfrac{7}{\sqrt{x}}\,dx$}
\q{CXVII}{Evaluate: $\displaystyle\int(5x^{2}+7x-3)^{7}(10x+7)\,dx$}
\q{CXVIII}{Evaluate: $\displaystyle\int\left(x+\dfrac{1}{x}\right)\left(1-\dfrac{1}{x^{2}}\right)dx$}
\q{CXIX}{Evaluate: $\displaystyle\int\tan x\,dx$}
\q{CXX}{Evaluate: $\displaystyle\int\sec x\,dx$}
\q{CXXI}{Evaluate: $\displaystyle\int\tan^{2}x\,dx$}
\q{CXXII}{Evaluate: $\displaystyle\int\cos^{2}x\,dx$}
\q{CXXIII}{Evaluate: $\displaystyle\int\sin 3x\cos 7x\,dx$}
\q{CXXIV}{Evaluate the following integrals: $\displaystyle\int\dfrac{ax+b}{\sqrt{ax^{2}+2bx+c}}$}
\q{CXXV}{Evaluate the following integrals: $\displaystyle\int x\sqrt{1+\sqrt{x}}\;dx$}
\q{CXXVI}{Evaluate the following integrals: $\displaystyle\int\dfrac{x}{1+x\tan x}\,dx$}
\q{CXXVII}{Evaluate the following integrals: $\displaystyle\int\dfrac{x}{(a+bx)^{3}}\,dx$}
\q{CXXVIII}{Evaluate the following integrals: $\displaystyle\int\dfrac{\sec^{2}x}{\tan x\,.\,\ln\tan x}\,dx$}
\q{CXXIX}{Evaluate the following integrals: $\displaystyle\int\dfrac{\ln\sin x}{\tan x}\,dx$}
\q{CXXX}{Evaluate the following integrals: $\displaystyle\int\dfrac{e^{2x}}{(2+e^{2x})^{2}}\,dx$}
\q{CXXXI}{Evaluate the following integrals: $\displaystyle\int\dfrac{e^{2x}-1}{e^{2x}+1}\,dx$}
\q{CXXXII}{valuate the following integral: $\displaystyle\int\dfrac{e^{x}}{9-e^{2x}}\,dx$}
\q{CXXXIII}{Evaluate the following integral: $\displaystyle\int\dfrac{e^{x}}{1+e^{2x}}\,dx$}
\q{CXXXIV}{Evaluate the following integral: $\displaystyle\int\dfrac{1}{x\sqrt{4-(\ln x)^{2}}}\,dx$}
\q{CXXXV}{Evaluate the following integral: $\displaystyle\int\dfrac{1}{9x^{2}-25}\,dx$}
\q{CXXXVI}{Evaluate the following integral: $\displaystyle\int\dfrac{1}{x\sqrt{3x^{2}-12}}\,dx$}
\q{CXXXVII}{Evaluate: $\displaystyle\int(5x^{5}+5)^{3}x^{4}\,dx$}
\q{CXXXVIII}{Evaluate: $\displaystyle\int\dfrac{\sin x}{e^{\cos x}}\,dx$}
\q{CXXXIX}{Evaluate: $\displaystyle\int\dfrac{e^{\tan^{-1}x}}{1+x^{2}}\,dx$}
\q{CXL}{Evaluate: $\displaystyle\int\dfrac{e^{x}\ln(1+e^{x})}{1+e^{x}}\,dx$}
\q{CXLI}{Evaluate: $\displaystyle\int\dfrac{(x+1)e^{x}}{1+x^{2}e^{2x}}\,dx$}
\q{CXLII}{Evaluate: $\displaystyle\int\dfrac{2x}{\sqrt{x^{4}+9}}\,dx$}
\q{CXLIII}{Evaluate the following integrals: Evaluate: $\displaystyle\int\ln x\,dx$}
\q{CXLIV}{Evaluate the following integrals: Evaluate: $\displaystyle\int x^{7}\,.\,\ln x\,dx$}
\q{CXLV}{Evaluate the following integrals: Evaluate: $\displaystyle\int x\cos 2x\,dx$}
\q{CXLVI}{Evaluate the following integrals: Evaluate: $\displaystyle\int\dfrac{x}{\sin^{2}x}\,dx$}
\q{CXLVII}{Evaluate the following integrals: Evaluate: $\displaystyle\int xe^{2x}\,dx$}
\q{CXLVIII}{Evaluate the following integrals: Evaluate: $\displaystyle\int x^{2}e^{ax}\,dx$}
\q{CXLIX}{Evaluate the following integrals: Evaluate: $\displaystyle\int x\tan^{2}x\,dx$}
\q{CL}{Evaluate the following integrals: Evaluate: $\displaystyle\int\ln(\sin x)\,.\cos x\,dx$}
\q{CLI}{Evaluate the following integrals: Evaluate: $\displaystyle\int e^{x}\ln(1+e^{x})\,dx$}
\q{CLII}{Evaluate the following integrals: Evaluate: $\displaystyle\int e^{x}\tan^{-1}(e^{x})\,dx$}
\q{CLIII}{Evaluate the following integrals: Evaluate: $\displaystyle\int\sqrt{x}\,\sin\sqrt{x}\,dx$}
\q{CLIV}{Evaluate the following integral: $\displaystyle\int\dfrac{e^{\tan^{-1}x}}{\sqrt{x}}\,dx$}
\q{CLV}{Evaluate the following integrals: Evaluate: $\displaystyle\int\sin^{-1}x\,dx$}
\q{CLVI}{Evaluate $\displaystyle\int x\sin nx\,dx,\;n\neq0$}
\q{CLVII}{Evaluate $\displaystyle\int xe^{ax}\,dx,\;a\neq0$}
\q{CLVIII}{Evaluate $\displaystyle\int\dfrac{x}{\cos^{2}x}\,dx$}
\q{CLIX}{Evaluate $\displaystyle\int e^{x^{2}}(\cos x+2\sin x)\,dx$}
\q{CLX}{Evaluate the following definite integral: $\displaystyle\int_{1}^{2}(x^{2}+x)\,dx$}
\q{CLXI}{Evaluate the following definite integral: $\displaystyle\int_{-2}^{3}(4x^{3}+3x^{2}+2x)\,dx$}
\q{CLXII}{Evaluate the following definite integral: $\displaystyle\int_{0}^{\pi/2}\sin^{2}x\,dx$}
\q{CLXIII}{Evaluate the following definite integral: $\displaystyle\int_{0}^{\pi/4}\dfrac{1}{1+\sin x}\,dx$}
\q{CLXIV}{Evaluate the following definite integral: $\displaystyle\int_{0}^{\pi/2}\cos^{3}\theta\,d\theta$}
\q{CLXV}{Evaluate the following definite integral: $\displaystyle\int_{0}^{1}\dfrac{\sqrt{x}}{1+x}\,dx$}
\q{CLXVI}{Evaluate the following definite integral: $\displaystyle\int_{0}^{1}\tan^{-1}x\,dx$}
\q{CLXVII}{Evaluate the following definite integral: $\displaystyle\int_{-3}^{5}f(x)\,dx$, where $f(x)=|x-4|+|x|$}
\q{CLXVIII}{Evaluate the following definite integral: $\displaystyle\int_{-3}^{4}|4-x^{2}|\,dx$}
\q{CLXIX}{Simplify: $\displaystyle\int_{0}^{1}f(x)\,dx-\int_{0}^{2}f(x)\,dx-\int_{3}^{1}f(x)\,dx+\int_{3}^{2}f(x)\,dx$}
\q{CLXX}{Evaluate the following definit integral: $\displaystyle\int_{0}^{4}x\sqrt{4-x}\,dx$}
\q{CLXXI}{Evaluate the following definit integral: $\displaystyle\int_{-1}^{1}\dfrac{e^{x}-1}{e^{x}+1}\,dx$}
\q{CLXXII}{Given the function f(x) satisfies the following condition: $\displaystyle\int_{0}^{4}f(x)\,dx=12$ evaluate $\displaystyle\int_{-1}^{5}f(x)\,dx$.}
\q{CLXXIII}{Show that: $\displaystyle\int_{\pi/6}^{\pi/3}\ln(\tan x)\,dx=0$}
\q{CLXXIV}{Evaluate: $\displaystyle\int_{-1}^{1}\left(x^{\frac{1}{5}}+1\right)dx$}
\q{CLXXV}{Evaluate: $\displaystyle\int_{0}^{\pi/4}\dfrac{\sec\theta}{\sin\theta+\cos\theta}\,d\theta$}
\q{CLXXVI}{Evaluate: $\displaystyle\int_{0}^{2}2xe^{x^{2}}\,dx$}
\q{CLXXVII}{Evaluate $\displaystyle\int_{0}^{2\pi}\cos^{5}x\,dx$}
\q{CLXXVIII}{Find the area of the region that i above the x-axis, but below the curve y=(x-2)(3-x)}
\q{CLXXIX}{Find the area of the region bounded by $y=|x-3|,x=2,\ x=4$ and the x-axis.}
\q{CLXXX}{Find the area bounded by the parabola y=x$^{2}$ -6x+8,x=2,x=4}
\q{CLXXXI}{For the function f(x)=2x over the interval [-3,3]: Find the net signed area}
\q{CLXXXII}{For the function f(x)=2x over the interval [-3,3]: Find the total (physical) area}
\q{CLXXXIII}{Find the moment of inertia of a uniform triangular lamina about its base.}
\q{CLXXXIV}{Find the integer b such that the average value of $f(x)=2+4x-3x^{2}$ on the interval [2,b] equals -7.}
\q{CLXXXV}{Find the average value of the function $f(x)=x^{2}+2$ over the interval [-1,2].}
\q{CLXXXVI}{A patient is given an intravenous dose oof an antibiotic. The concentration in the bloodstream over time is modeled by: $C(t)=8e^{-0.15t}(mg/L)$ where t is an hours. What is the AUC from t=0 to infinity?}
\q{CLXXXVII}{A patient is given an intravenous dose oof an antibiotic. The concentration in the bloodstream over time is modeled by: $C(t)=8e^{-0.15t}(mg/L)$ where t is an hours. If the minimum required total exposure over the first 6 hours is 40 mg. h/L, is the drug exposure from thid dose sufficient for a 6-hour treatment?}
\q{CLXXXVIII}{Write the order and degree (if ddefined) of the following differential equation: $\dfrac{d^{3}y}{dx^{3}}+7\dfrac{d^{2}y}{dx^{2}}-9xy\dfrac{dy}{dx}+x^{2}y=0$}
\q{CLXXXIX}{Write the order and degree (if ddefined) of the following differential equation: $\left(\dfrac{d^{2}y}{dx^{2}}\right)^{3}+\dfrac{dy}{dx}=1$}
\q{CXC}{Write the order and degree (if ddefined) of the following differential equation: $\dfrac{d^{2}y}{dx^{2}}+\sin\!\left(\dfrac{dy}{dx}\right)=0$}
\q{CXCI}{Write the order and degree (if defined) of the following differential equation: $\dfrac{dy}{dx}+e^{y}\sin(x)=0$}
\q{CXCII}{Write the order and degree (if defined) of the following differential equation: $\cos x\left(\dfrac{d^{2}y}{dx^{2}}\right)^{\frac{1}{3}}+\dfrac{dy}{dx}+y=0$}
\q{CXCIII}{Show that the following relation written against the differential equation is its solution. $x\dfrac{dy}{dx}=1+y\ ;\ y=cx-1$}
\q{CXCIV}{Show that the following relation written against the differential equation is its solution. $x^{2}(2y+1)\dfrac{dy}{dx}=1\ ;\ y^{2}+y=5-\dfrac{1}{x}$}
\q{CXCV}{Show that the following relation written against the differential equation is its solution. $\dfrac{ds}{dt}=e^{t}(s^{2}+1)\ ;\ s=\tan(e^{t}-10)$}
\q{CXCVI}{Is the folowing equation is separable or homogeneous? $x\dfrac{dy}{dx}=1+y$}
\q{CXCVII}{Is the folowing equation is separable or homogeneous? $x^{2}(2y+1)dy-dx=0$}
\q{CXCVIII}{Is the folowing equation is separable or homogeneous? $(2xy+x^{2})dy-y^{2}dx=0$}
\q{CXCIX}{Is the folowing equation is separable or homogeneous? $\dfrac{dy}{dx}=e^{x-y}$}
\q{CC}{Is the folowing equation is separable or homogeneous? $\dfrac{dT}{dt}+k(T-T_{m})=0$}
\q{CCI}{Is the folowing equation is separable or homogeneous? $\dfrac{dy}{dx}=\ln(x)-\ln(y)\ ;\ x>0,\ y>0$}
\q{CCII}{Find the general solution of the differential equation: $\dfrac{dy}{dx}=\left(\dfrac{1-y^{2}}{1-x^{2}}\right)^{\frac{1}{2}}$}
\q{CCIII}{Find the general solution of the differential equation: $\dfrac{dy}{dx}=e^{x-y}$}
\q{CCIV}{Find the general solution of the differential equation: $\dfrac{dy}{dx}=2x\sqrt{1-y^{2}}$}
\q{CCV}{Find the general solution of the differential equation: $x^{2}dx+(1-x^{2})dy=0$}
\q{CCVI}{Find the general solution of the differential equation: $x\sin x\,dx+\ln y\,dy=0$}
\q{CCVII}{Find the general solution of the differential equation: $(y^{2}+2xy)dx+x^{2}dy=0$}
\q{CCVIII}{Find the general solution of the differential equation: $(x^{2}+3y^{2})dx-2xy\,dy=0$}
\q{CCIX}{Find the general solution of the differential equation: $(x^{2}+xy+y^{2})dx-x^{2}dy=0$}
\q{CCX}{Find the general solution of the differential equation: $x\dfrac{dy}{dx}=xe^{-\frac{y}{x}}+y$}
\q{CCXI}{Find the general solution of the differential equation: $y-x\dfrac{dy}{dx}=3\left(1+x\dfrac{dy}{dx}\right)$}
\q{CCXII}{Find the general solution of the differential equation: $y\,dx+x\,dy=0$}
\q{CCXIII}{Solve the initial value problem: $\dfrac{dy}{dx}-x=xy^{2}\ ;\ y(0)=1$}
\q{CCXIV}{Solve the initial value problem: $\dfrac{ds}{dt}+2st=0\ ;\ s=4e$ when $t=0$}
\q{CCXV}{Solve the initial value problem: $\dfrac{dy}{dx}=\dfrac{x+y}{x}\ ;\ y(1)=1$}
\q{CCXVI}{Show that the relation $y=cx+2\ldots.(1)$ represents a solution of the differential equation $y-x\dfrac{dy}{dx}=2$}
\q{CCXVII}{Find the general solution of the differential equation $x\,dx+y\,dy=0$}
\q{CCXVIII}{Solve the initial value problem: $\dfrac{dy}{dx}=\dfrac{3}{4}x^{2}+x-3\ ;\ y(2)=0$}
\q{CCXIX}{A particle is moving in a straight line and its acceleration is given by a=2t-7 Find: velocity v in terms of if v=10 m/sec, when t=0}
\q{CCXX}{A particle is moving in a straight line and its acceleration is given by a=2t-7 Find: Distance s in terms of t if s=0 m when r=0.}
\q{CCXXI}{A radioactive isotope has a half-life of 12 hours. The initial mass is 160 grams. What is the decay constant?}
\q{CCXXII}{Find equation of the altitude of triangle XYZ ffrom vertex X(1,3), Y(5,7) and Z(-3,9).}
\q{CCXXIII}{Show that the medians of the triangle with vertices A(3,2), B(-1,-5) and C(7,2) are concurrent.}
\q{CCXXIV}{Convert the following equation into: 4x+5y-20=0 (i) Slope intercept form (ii) Two point form (iii) Intercept form (iv) Point slope form (v) Normal form (vi) Symmetric form.}
\q{CCXXV}{Find the value of k if the lines x+y-4=0, 2x-y+1=0 and kx+2y-10=0 are cocurrent.}
\q{CCXXVI}{If the lines 2x+3y-7=0, 4x-y-5=0 an 6x+ay-9=0 meet at a point then find the value of a.}
\q{CCXXVII}{If the lines 2x+3y-7=0, 4x-y-5=0 an 6x+ay-9=0 meet at a point then find the value of a.}
\q{CCXXVIII}{If the lines 2x+3y-7=0, 4x-y-5=0 an 6x+ay-9=0 meet at a point then find the value of a.}
\q{CCXXIX}{Find an equation of the straight line if its slope is 3 and y-intercept is 7.}
\q{CCXXX}{Write down an equation of the straight line passing through (-4,3) and having slope -5.}
\q{CCXXXI}{Find an equation of the straight line passing through the points (-3,-2) and (5,4).}
\q{CCXXXII}{Write down an equation of the straight line which cuts the x-axis at (3,0) and y-axis (0,-6).}
\q{CCXXXIII}{The length of perpendicular from the origin to a line is 3 units and the inclination of this perpendicular is $120^{o}$ . Find the equation of the straight line.}
\q{CCXXXIV}{Find the angle measured from the line $\ell_{1}$ to the line $\ell_{2}$, where. $\ell_{1}=$ joining $(2,7)$ and $(-4,5)$. Also find the acute angle.}
\q{CCXXXV}{Find the angle measured from the line $\ell_{1}$ to the line $\ell_{2}$, where. $\ell_{1}=$ joining $(1,2)$ and $(3,5)$. Also find the acute angle.}
\q{CCXXXVI}{Find the angle from the line with slope m$_{1}$ to the line with slope m$_{2}$ , where $m_{1}=2,\ m_{2}=-\dfrac{1}{2}$}
\q{CCXXXVII}{Find the angle from the line with slope m$_{1}$ to the line with slope m$_{2}$ , where $m_{1}=3,\ m_{2}=-2$}
\q{CCXXXVIII}{Find the angle from line $\ell_{1}$ to the line $\ell_{2}$, where: $\ell_{1}$:$2x+y-3=0$ ; $\ell_{2}$: $x+y-5=0$}
\q{CCXXXIX}{Find the angle from line $\ell_{1}$ to the line $\ell_{2}$, where: $\ell_{1}$:$2x-3y+2=0$ ; $\ell_{2}$: $x+y-4=0$}
\q{CCXL}{Find the angle from the line with slope $-\dfrac{7}{3}$ to the line with slope $\dfrac{5}{2}$.}
\q{CCXLI}{Find the area of the region bounded by the triangle with vertices A(a,b), B(-a,b) and C(0,-b).}
\q{CCXLII}{Find the area of the region bounded by the triangle with vertices A(1, 2), B(3, 6) and C(5, 10).}
\q{CCXLIII}{Find the equation of lines represented by the following and also find measure of the angle between them: $6x^{2}-5xy-4y^{2}=0$}
\q{CCXLIV}{Find the equation of lines represented by the following and also find measure of the angle between them: $2x^{2}-7xy+3y^{2}=0$}
\q{CCXLV}{Find the equation of lines represented by the following and also find measure of the angle between them: $x^{2}+2xy\sec\alpha+y^{2}=0$}
\q{CCXLVI}{A(-2,-1) and B(3,4) are two points on the solar panel roof. Find slope and tilt angle.}
\q{CCXLVII}{Find an equation of each of lines represented by $3x^{2}-10xy+3y^{2}=0$}
\q{CCXLVIII}{Find measure of the acute angle between the lines represented by: $5x^{2}+6xy+y^{2}=0$}
\q{CCXLIX}{Planet A is at (1,0) AU and planet B is at (-5,0 0.866) AU. Find the distance between them. Write answer in km.}
\q{CCL}{Air craft flies from A(100,200) km to B(500,800) km. Find slope of the track and straight line distance.}
\q{CCLI}{In the following find an equation of the circle with: centre at (5,7) and radius 2.}
\q{CCLII}{In the following find and equation of the circle with centre at $(\alpha\cos\theta,\ \alpha\sin\theta)$ and radius a.}
\q{CCLIII}{In the following find and equation of the circle with ends of diameter (3,4) and (2,-7).}
\q{CCLIV}{Express the following in the standard form and hence find its centre and radius: $2x^{2}+2y^{2}-4x-12y+11=0$}
\q{CCLV}{Express the following in the standard form and hence find its centre and radius: $x^{2}+y^{2}-2ax+2by-2ab=0$.}
\q{CCLVI}{Find an equation of cricle passing through the points (2,8), (5,-1) and (-3,3) also find the centre and radius.}
\q{CCLVII}{Find an equation of a cirlce of radius a lying in the second quadrant and tangent to both the axes.}
\q{CCLVIII}{Find the equation of the circle which touches the x-axis at the point (5,0) and passed through another point (3,2).}
\q{CCLIX}{Find the equations of the circles of radius 5 that are tangent to the line 4x+3y-10=0 at the point (1,2).}
\q{CCLX}{Write an equation of the circle with centre (-2,1) and radius 5.}
\q{CCLXI}{Find the equations of the tangent and normal to the circles: $x^{2}+y^{2}=8$ at $(2,2)$ and at $(2\sqrt{2}\cos\theta,\ 2\sqrt{2}\sin\theta)$}
\q{CCLXII}{Find the equations of the tangent and normal to the circles: $x^{2}+y^{2}-6x-2y-7=0$ at $(4,5)$.}
\q{CCLXIII}{Show that the line 4x+3y-25=0 is tangent to the circle $x^{2}+y^{2}=25$ and find the point of contact.}
\q{CCLXIV}{Find the equation of the tangents to the circle $x^{2}+y^{2}=4$: Which are perpendicular to the line 6x-8y+3=0.}
\q{CCLXV}{Determine whether the point P(2,3) lies outside. on, or inside the circle: $x^{2}+y^{2}+2x-3y-10=0$.}
\q{CCLXVI}{Find the length of the tangent fro point P(2,-5) to the circle $5x^{2}+5y^{2}+10x+14y-15=0$.}
\q{CCLXVII}{Find the focus, vertex and directrix of the following parabola: $y^{2}=12x$.}
\q{CCLXVIII}{Find the focus of the following parabola: $x^{2}=8(y-1)$.}
\q{CCLXIX}{Find the vertex of the following parabola: $x^{2}=8(y-1)$.}
\q{CCLXX}{Find the directrix of the following parabola: $x^{2}=8(y-1)$.}
\q{CCLXXI}{Find the focus of the following parabola: $y^{2}=4(x-1)$.}
\q{CCLXXII}{Find the vertex of the following parabola: $y^{2}=4(x-1)$.}
\q{CCLXXIII}{Find the directrix of the following parabola: $y^{2}=4(x-1)$.}
\q{CCLXXIV}{Write the equation of the parabola with given elements: Focus (3,5); vertex (3,2).}
\q{CCLXXV}{Write the equation of the parabola with given elements: Directrix x=4, vertex (6,3).}
\q{CCLXXVI}{Write the equation of the parabola with given elements: Axis y=0, through (2,2) and (6,-3).}
\q{CCLXXVII}{Find the equation of the parabola gaving its focus at the origin and directrix parallel to: (i) x-axis}
\q{CCLXXVIII}{Find the equation of the parabola gaving its focus at the origin and directrix parallel to: (i) y-axis.}
\q{CCLXXIX}{Prove that in a parabola, a circle described on the latus rectum as diameter touches the directrix.}
\q{CCLXXX}{Find the equation of the tangent and normal to the parabola $y^{2}=4ax$ at the point $(at^{2},2at)$.}
\q{CCLXXXI}{Find an equation of the ellipse with given data: Foci $(0,\pm5)$ and minor axis of length 8.}
\q{CCLXXXII}{Find an equation of the ellipse with given data: Vertices (-2,5) and (-2,-5); foci (-2,4) and (-2,-2).}
\q{CCLXXXIII}{Find an equation of the ellipse with given data: Vertices $(\pm6,0)$, eccentricity $\dfrac{2}{3}$}
\q{CCLXXXIV}{Find the equation of an ellipse centered at (5,3) that is tangent to both te x-axis and y-axis.}
\q{CCLXXXV}{Find the minimum distance of the point P($3\cos\theta,\ 2\sin\theta$) on the ellipse from the centre of the ellipse.}
\q{CCLXXXVI}{Find the minimum distance from the ellipse $\dfrac{x^{2}}{9}+\dfrac{y^{2}}{4}=1$ to the line x+y=6.}
\q{CCLXXXVII}{An ellipse has the following properties: Its major axis is vertical and has a length of 10 units.}
\q{CCLXXXVIII}{An ellipse has the following properties: The distance between its foci is 6 units.}
\q{CCLXXXIX}{An ellipse has the following properties: The centre is at the point (3,-2). Find the standard equation of this ellipse.}
\q{CCXC}{Find an equation of the hyperbola with the given data: centre (0,0), Focus (0,10), Vertex (0,8).}
\q{CCXCI}{Find the equation of the hyperbola with the given data: Foci $(0,\pm8)$, Eccentricity=4}
\q{CCXCII}{Find the equation of the hyperbola with the given data: Foci $(0,\pm10)$, Directrices $y=\pm5$.}
\q{CCXCIII}{Find the centre, foci, eccentricity, vertices and equations of directrices of the following: $\dfrac{x^{2}}{36}-\dfrac{y^{2}}{64}=1$}
\q{CCXCIV}{Find the centre, foci, eccentricity, vertices and equations of directrices of the following: $\dfrac{(x+3)^{2}}{5}-\dfrac{(y-2)^{2}}{4}=1$}
\q{CCXCV}{Find the centre, foci, eccentricity, vertices and equations of directrices of the following: $\dfrac{y^{2}}{25}-\dfrac{x^{2}}{16}=1$}
\q{CCXCVI}{Find the equation of the hyperbola in standard form (centred at the origin) that passes through the points $(\sqrt{80},4)$ and $(\sqrt{32},2)$.}
\q{CCXCVII}{Find equations of the tangent and normal to the hyperbola $\dfrac{x^{2}}{a^{2}}-\dfrac{y^{2}}{b^{2}}=1$ at the point $(a\sec\theta,\ b\tan\theta)$.}
\q{CCXCVIII}{Find the length of the latus rectum of the hyperbola $9x^{2}-16y^{2}=144$.}
\q{CCXCIX}{Find the eccentricity of the hyperbola whose asymptotes are perpendicular.}
\q{CCC}{If the transverse axis of a hyperbola is 8 and the latus rectum is 9, find its eccentricity.}
\q{CCCI}{Find the eccentriity of the hyperbola whose latus rctum is half its tranverse axis.}
\q{CCCII}{If the distance between the foci of a hyperbola is 16 and the distance between its vertices is 12, find its eccentricity.}
\q{CCCIII}{Find the equation of the hyperbola whose asymptotes are $y=\pm2x$ and which passes through the point (3,4).}
\q{CCCIV}{A parabolic arch has 120 m base and height 25 m. Find the height of the arch at a point 42 m from centre of the base.}
\q{CCCV}{A parabolic flaslight reflector has the shape given by the equation: $x=\dfrac{1}{20}y^{2}$ where should the lightbulb be placed so that the reflected rays emerge parallel to the axis?}
\q{CCCVI}{A vehicle moving with a speed of 36 km/h is stopped by applying brakes and comes to rest in 10 seconds. Find the retardation and distance covered during this time.}
\q{CCCVII}{The acceleration of a particle is $a(t)=\dfrac{1}{1+t}$ for $t\geq0$. If the velocity is 0 m/s when t=0, find the velocity at t=5 seconds.}
\q{CCCVIII}{A man walks 100 meters east in 2 minutes, then turns around and walks 60 meters west 1 minute. Find his speed and velocity.}
\q{CCCIX}{A car starts from rest and accelerates uniformyl at 3 m/s$^{2}$ Find the velocity of the car after 5 seconds.}
\q{CCCX}{A car accelerates from 10 m/s to 30 m/s in 4 seconds. Find the distance covered during this time.}
\q{CCCXI}{Find the position and the velocity of a particle moving along a straight line at any time t seconds, given that it starts rest at t=0 and is subject to an acceleration a(t)=t$^{2}$ +cost}
\q{CCCXII}{The acceleration (in m/s$^{2}$) of a particle is given a(t)=xost+sint.}
\q{CCCXIII}{Use Trapezoidal Rule to approximate the value of the following definite integrals. $\displaystyle\int_{0}^{2}\cos(x)\,dx$ taking n=4}
\q{CCCXIV}{Find the principal value of the following: $\sin^{-1}\!\left(\dfrac{\sqrt{3}}{2}\right)$}
\q{CCCXV}{Find the principal value of the following: $\cos^{-1}\!\left(-\dfrac{\sqrt{2}}{2}\right)$}
\q{CCCXVI}{Find the principal value of the following: $\sec^{-1}\!\left(\dfrac{2}{\sqrt{3}}\right)$}
\q{CCCXVII}{Find the principal value of the following: $\csc^{-1}(-\sqrt{2})$}
\q{CCCXVIII}{Find the principal value of the following: $\sin^{-1}\!\left[\cos\!\left(\dfrac{3\pi}{4}\right)\right]$}
\q{CCCXIX}{Find the principal value of the following: $\tan^{-1}\!\left[\cot\!\left(\dfrac{\pi}{6}\right)\right]$}
\q{CCCXX}{Find the principal value of the following: $\sec^{-1}\!\left[2\cos\!\left(\dfrac{3\pi}{4}\right)\right]$}
\q{CCCXXI}{Find the principal value of the following: $\csc^{-1}\!\left(\sqrt{2}\cos\dfrac{\pi}{4}\right)$}
\q{CCCXXII}{Find the domain of the following: $f(x)=\sin^{-1}(5x-3)$}
\q{CCCXXIII}{Find the domain of the following: $g(x)=\cos^{-1}\sqrt{x-3}$}
\q{CCCXXIV}{Find the domain of the following: $h(x)=\csc^{-1}(x^{2}-3)$}
\q{CCCXXV}{Find the domain of the following: $k(x)=\sin^{-1}(\ln x)$}
\q{CCCXXVI}{Find the domain of the following: $p(x)=\tan^{-1}\sqrt{4-x^{2}}$}
\q{CCCXXVII}{Find the domain of the following: $q(x)=\sec^{-1}(x+2)$}
\q{CCCXXVIII}{Find the domain of the following: $u(x)=\sec^{-1}\!\left(\dfrac{3x-1}{2}\right)$}
\q{CCCXXIX}{Find the domain of the following: $v(x)=\sin^{-1}(\ln x)+\csc^{-1}(e^{x})$}
\q{CCCXXX}{Without using calculator find the value of the following: $\cos\!\left(\dfrac{1}{2}\cot^{-1}\dfrac{5}{12}\right)$}
\q{CCCXXXI}{Without using calculator find the value of the following: $\sec\!\left(\dfrac{1}{2}\sin^{-1}\dfrac{7}{25}\right)$}
\q{CCCXXXII}{Without using calculator find the value of the following: $\tan\!\left(\dfrac{1}{2}\sin^{-1}\dfrac{\sqrt{3}}{2}\right)$}
\q{CCCXXXIII}{Find the principal value of: $\sin^{-1}\!\left(\dfrac{1}{2}\right)$}
\q{CCCXXXIV}{Find the principal value of: $\cos^{-1}\!\left(-\dfrac{\sqrt{3}}{2}\right)$}
\q{CCCXXXV}{Find the principal value of: $\tan^{-1}\!\left(-\dfrac{1}{\sqrt{3}}\right)$}
\q{CCCXXXVI}{Find value of: $\sec\left[\tan^{-1}(-\sqrt{3})\right]$}
\q{CCCXXXVII}{Find value of: $\cot\left[\cos^{-1}\!\left(-\dfrac{1}{2}\right)\right]$}
\q{CCCXXXVIII}{Find the domain of the following functions: $f(x)=\cos^{-1}(3x-1)$}
\q{CCCXXXIX}{Find the domain of the following functions: $g(x)=\sin^{-1}\!\left(\dfrac{x^{2}-5}{4}\right)$}
\q{CCCXL}{Find the domain of the following functions: $h(x)=\sec^{-1}(\ln x)$}
\q{CCCXLI}{Prove that: $\csc^{-1}x=\sin^{-1}\dfrac{1}{x},\ x\in(-\infty,-1]\cup[1,\infty)$}
\q{CCCXLII}{Prove that: $\sec^{-1}x=\cos^{-1}\dfrac{1}{x},\ x\in(-\infty,-1]\cup[1,\infty)$}
\q{CCCXLIII}{Prove that: $\cot^{-1}x=\tan^{-1}\dfrac{1}{x},\ x\in(0,\infty)$}
\q{CCCXLIV}{Without using calculator find the value of $\sin\left[\dfrac{1}{2}\tan^{-1}\dfrac{4}{3}\right]$.}
\q{CCCXLV}{Prove the following: $2\tan^{-1}x=\cos^{-1}\!\left(\dfrac{1-x^{2}}{1+x^{2}}\right)$}
\q{CCCXLVI}{Prove the following: $\cos(\sin^{-1}x)=\sqrt{1-x^{2}}$}
\q{CCCXLVII}{Prove the following: $\cot^{-1}\!\left(\dfrac{\cos x}{1-\sin x}\right)=\dfrac{\pi}{4}-\dfrac{x}{2}$}
\q{CCCXLVIII}{Prove that $\cos^{-1}x=\sin^{-1}\sqrt{1-x^{2}},\ x\in[0,1]$}
\q{CCCXLIX}{Solve the following trigonometric eqautions: $\sin x=\dfrac{1}{\sqrt{2}}$}
\q{CCCL}{Solve the following trigonometric eqautions: $2\cos x=-1$}
\q{CCCLI}{Solve the following trigonometric eqautions: $3\tan^{2}x=1$}
\q{CCCLII}{Solve the following trigonometric equations: $9-3\csc^{2}x=5$}
\q{CCCLIII}{Solve the following trigonometric eqautions: $2\cos^{2}x-3\cos x+1=0$}
\q{CCCLIV}{Solve the following trigonometric eqautions: $\sin^{2}x-3\cos^{2}x=0$}
\q{CCCLV}{Solve the following trigonometric equations: $\sqrt{2}\cos x\sin x-\sin x=0$}
\q{CCCLVI}{Solve the following trigonometric eqautions: $3\cot^{2}x+3\csc^{2}x=5$}
\q{CCCLVII}{Solve the following trigonometric eqautions: $\sin 2x-\sqrt{3}\cos x=0$}
\q{CCCLVIII}{Solve the following trigonometric eqautions: $\cos 2x=\cos x$}
\q{CCCLIX}{Solve the following trigonometric eqautions: $\tan x=\sin 2x$}
\q{CCCLX}{Solve the following trigonometric eqautions: $\sin 3x=\cos 2x$}
\q{CCCLXI}{Solve $\sqrt{3}\cot x-\csc x-1=0$}
\q{CCCLXII}{Solve graphically: $3\sin x=\dfrac{x}{2}+1,\ x\in[-\pi,\pi]$}
\q{CCCLXIII}{Solve graphically: $\cos(2x-1)=\dfrac{x}{5},\ x\in[0,\pi]$}
\q{CCCLXIV}{Give two real-life examples where vector-valued functions are necessary.}
\q{CCCLXV}{Construct a vector function for a particle moving in a circle.}
\q{CCCLXVI}{Write a vector function for a car moving along a parabolic path.}
\q{CCCLXVII}{Create a 3D vector function for a drone flying upward in a spiral.}
\q{CCCLXVIII}{Given $x(t)=t^{2}+1,y(t)=\sqrt{t}$ construct, construct r(t).}
\q{CCCLXIX}{Find the Find the domain and range of the following vector-valued functions: $r(t)=\sqrt{t}\,\hat{i}+\ln t\,\hat{j}$}
\q{CCCLXX}{Find the Find the domain and range of the following vector-valued functions: $r(t)=\cos(t)\,\hat{i}+\sin(t)\,\hat{j}$}
\q{CCCLXXI}{Find the Find the domain and range of the following vector-valued functions: $r(t)=t\hat{i}+t^{2}\hat{j}+e^{p}\hat{k}$}
\q{CCCLXXII}{Find the first-order derivatives of the following vector-valued functions with respect to t. $\boldsymbol{r}(t)=(t^{2}+1)\hat{i}+(t^{3}+t^{2}+3t)\hat{j}$}
\q{CCCLXXIII}{Find the first-order derivatives of the following vector-valued functions with respect to t. $\boldsymbol{r}(t)=(t\sin t-\cos t)\,\hat{i}+(\cot(t)+\ln t)\hat{j}+(\ln(1+t^{2}))\hat{k}$}
\q{CCCLXXIV}{Find the first derivative of $r(t)=t^{2}\hat{i}+\sin t\,\hat{j}+e^{t}\hat{k}$ with respect to t.}
\end{document}
